\begin{table}[htbp]
\centering
\small
\caption{V2 H1--H2（幾何再編と表現共有度の変位）：事後学習に伴う表現幾何学的歪み（Procrustes Distortion）、デコードピーク深度変位 $\Delta d^*$、およびタスク共有度変化 $\Delta\text{Sharing}$。4ファミリー統合ブートストラップ95\%信頼区間。}
\label{tab:v2_h1_h2_reorganization}
\begin{tabular}{ll ccc}
\toprule
\textbf{Hypothesis} & \textbf{Metric} & \textbf{Estimate} & \textbf{95\% CI} & \textbf{Condition} \\
\midrule
\multirow{4}{*}{H1a: Geometric Distortion} & Reader Procrustes Distortion         & ---      & ---                & Matched-Plain \\
                                 & Self Procrustes Distortion           & ---      & ---                & Matched-Plain \\
                                 & RSA Reader ($\rho_{\text{RSA}}$)     & ---      & ---                & Matched-Plain \\
                                 & RSA Self ($\rho_{\text{RSA}}$)       & ---      & ---                & Matched-Plain \\
\midrule
\multirow{4}{*}{H1b: Decodability Peak Shift} & Valence Reader $\Delta d^*$          & ---      & ---                & Matched-Plain \\
                                 & Valence Self $\Delta d^*$            & ---      & ---                & Matched-Plain \\
                                 & Arousal Reader $\Delta d^*$          & ---      & ---                & Matched-Plain \\
                                 & Arousal Self $\Delta d^*$            & ---      & ---                & Matched-Plain \\
\midrule
\multirow{2}{*}{H2: Sharing Reorganization} & $\Delta\text{Sharing}$ (Valence)     & ---      & ---                & Matched-Plain \\
                                 & $\Delta\text{Sharing}$ (Arousal)     & ---      & ---                & Matched-Plain \\
\bottomrule
\end{tabular}
\vspace{1ex}
\begin{minipage}{\linewidth}
\footnotesize
\textbf{Note:} $\Delta d^* = d^*_{\text{instruct}} - d^*_{\text{base}} > 0$ は、事後学習によって表現デコードピークが後段側（deeper側）へシフトしたことを示す。また、$\Delta\text{Sharing} < 0$ はReaderとSelfの表現共有度合いが事後学習によってタスク分離方向に再編されたことを示す。
\end{minipage}
\end{table}